\scanpage{30}
de se borner aux données de ramification $\underline R$ du type envisagé dans l'exemple 3; ce sont les seules sans doute qu'on rencontrera jamais, et les seules aussi où on sait construire le foncteur réalisation géométrique $\underline r$ et en dire des choses raisonnables.

Avec l'une ou l'autre approche mentionnée à l'alinéa précédent pour traiter de données de ramifications sur des schémas variables sur $S$, (et toujours sous les conditions de l'exemple 3 ...) on est en condition pour énoncer et démontrer qu\deleted{e}'un morphisme fidèlement plat et quasi-compact $T\quad S$ est de descente effective pour la catégorie fibrée sur $(S,T,T\mathbin{\mathrm{x}}_S T,T\mathbin{\mathrm{x}}_S T\mathbin{\mathrm{x}}_S T)$ formée des revêtements de type $\underline R_{(S)}$ ,..... Ceci s'appliquera en particulier à la \unsure{descente|desente} fidèlement plate des revêtements kummériens; c'est utilisé (sans démonstration) à la page 38 des notes de Murre dans le cas de la \unsure{desente|descente} à partir d'un anneau complété, sans compter diverses utilisations de descente à partir d'un localisé strict. Il y a également lieu d'énoncer une propriété de passage à la limite, pour un système projectif $(\unsure{S|X}_i)$ de schémas quasi-compacts et quasi-séparés à morphismes de transition affines; cela permettra de ramener ce qui se passe au voisinage d'un point $s$ d'un $S$ à voir ce qui se passe en $\mathrm{Spec}(\underline O_{S,s})$. Par exemple, dans le cas $S$ \inserted{loc noeth} normal munie d'une famille $(D_i)_{i\;\;I}$ de diviseurs réduits, donc où \inserted{la catégorie} des revêtements kummériens s'identifie\deleted{nt} à \deleted{\unsure{des}} une sous-catégorie de la catégorie des revêtements, on trouve que le revêtement $X$ de $S$ est \deleted{\unsure{norm}} kummérien sss pour tout $s\quad S$, il en est de même du revêtement induit sur $\mathrm{Spec}(\underline O_{S,s})$, et on peut dans cet énoncé remplacer $\underline O_{S,s}$ son hensélisé ou par son hensélisé strict, \deleted{Il y aura un énoncé \unsure{a}} voire par son complété si les hypothèses faites restent stables
% archive sheet number: 29
